\SourceLeaf{361}%
Ausführung der Rechnung bereits sehr gute Näherungswerte der zu bestimmenden Größen vorliegen. Genau genommen haben auch alle anderen Lösungsarten des Problems, insbesondere die von \textsc{Bessel}, teilweise den Charakter indirekter Rechenschritte. Das Wesentliche bleibt aber, dass von den wiederholten Annäherungen nur die letzte aufbewahrt werden muss. Sie enthält den gesamten Kern der Rechnung und ist so kurz und übersichtlich wie bei keiner anderen Methode.

Die Formeln zur Lösung der Aufgabe auf der Sphäroidoberfläche unterscheiden sich von denen für die Kugeloberfläche lediglich dadurch, dass gewisse Zwischengrößen, die bei diesen konstant sind. bei jenen von der Breite abhängen. Diese Größen lassen sich folglich in einer Hilfstafel zusammenstellen, deren Argument die Breite ist. Steht eine solche Hilfstafel zur Verfügung, so ist die Rechnung auf der Sphäroidoberfläche in jedem konkreten Fall genauso leicht wie auf der Kugeloberfläche. Für die Zone von 51---54 Grad Breite, die für das hannoversche Dreiecksnetz ausreicht, ist am Schluss der Abhandlung eine solche Hilfstafel beigefügt. Ihr liegt der Wert der Abplattung zugrunde, den \textsc{Bessel} aus allen bisherigen Gradmessungen abgeleitet hat und der schon der ersten Abhandlung zugrunde lag. Wer dieselbe Methode auf ein Dreiecksnetz außerhalb dieser Zone anwenden will, muss zunächst jene Hilfstafel für seinen Zweck weiter ausdehnen oder, falls er eine andere Abplattung zugrunde legen will, erst eine neue Hilfstafel erstellen. Bei der Bearbeitung eines großen Dreiecksnetzes fällt eine solche Vorarbeit überhaupt nicht ins Gewicht. Die Bequemlichkeit der Anwendung wiegt die dafür aufgewandte Mühe reichlich auf. Für den Fall dagegen, dass man nur eine oder ein paar konkrete Lösungen der betreffenden Aufgabe sucht, sollte die Methode nicht in erster Linie bestimmt sein.
